% Standalone paragraph. Paste into your own thesis, or \input it.
% Not wired into any chapter's numbering -- drop it wherever the "which features are actually
% reproducible" argument is needed. Companion to whole_tree_gap.tex: that file argues the field
% measures the wrong place in the tree, this one argues it does not agree on how to measure what
% it does look at.


\section*{Features}


One coronary feature is more standardized than the rest: bifurcation angle has a fixed definition,
because a stent has to be placed the same way every time, codified across eighteen consensus
documents from the European Bifurcation Club \cite{burzotta2024ebc}. I might want to look into
expanding it, though -- e.g.\ the angle could be the same whether a large vessel bifurcates into a
much thinner one or a thin vessel bifurcates into an equally thin one, so the same angle can hide
very different information.

Furthermore, I am going to focus much of my effort on developing a more mature method for
tortuosity. Right now its dominant metric, the tortuosity index, compares path
length to straight-line endpoint distance and, in the review's own words, ``cannot capture the
spatial information'' in the first place \cite{shen2026geometry}.
Clinical practice reduces it further still, to a bend count or a 2D C-/S-shape projection.
The formula itself, not just its use, is contested.

